%% ======================================================================
%%  Chapter V -- The splitting of prime ideals in an extension field
%% ======================================================================
\clearpage
\folio{68}%
\chapterhead{5}{V}{The splitting of prime ideals in an extension field}{THE SPLITTING OF PRIME IDEALS IN AN EXTENSION FIELD}
\renewcommand{\chaptertitle}{THE SPLITTING OF PRIME IDEALS IN AN EXTENSION FIELD}

Let $K$ be a number field, $A$ the ring of integers of $K$, $L$ an extension
of finite degree of $K$, and $B$ the integral closure of $A$ in $L$ (i.e.\
the ring of integers of $L$). The ideal $B\mathfrak{p}$, generated in $B$ by
a non-zero prime ideal $\mathfrak{p}$ of $A$, is not in general prime. It
splits into a product of prime ideals (\hyperref[sec:3.4]{Chapter III, §4}, Theorem
\ref{thm:III.4.3}), i.e.\
$B\mathfrak{p}=\prod_{i}\mathfrak{P}_{i}^{e_{i}}$. In this chapter we propose
to study this splitting. The case in which $B$ is a free $A$-module (for
example, when $A$ is a principal ideal ring; cf.\ \hyperref[sec:2.7]{Chapter II, §7}, the
corollary of Theorem \ref{thm:II.7.1}) is particularly easy. In \hyperref[sec:5.1]{§1} we shall
show how the general case may be reduced to this easier case.

%% ----------------------------------------------------------------------
\sechead{1}{Preliminaries concerning rings of fractions}%
        {PRELIMINARIES CONCERNING RINGS OF FRACTIONS}

\begin{definition}
\label{def:V.1.1}
Let $A$ be an integral domain, let $K$ be its field of fractions, and let $S$
be a subset of $A-(0)$ which is stable under multiplication and contains $1$.
We call the set of all elements of $K$ which may be written in the form $a/s$
with $a\in A$ and $s\in S$ the \emph{ring of fractions} of $A$ with respect
to $S$. We denote it $S^{-1}A$.
\end{definition}

Clearly $S^{-1}A$ is a commutative ring (since
$a/s+a'/s'=(s'a+sa')/ss'$ and $(a/s)\cdot(a'/s')=aa'/ss'$) which contains
$A$ (since $1\in S$). If $S=A-(0)$, then $S^{-1}A=K$. If $S$ contains $1$
alone, or if it contains only units of $A$, then $S^{-1}A=A$.

\begin{proposition}
\label{prop:V.1.1}
Let $A$ be an integral domain and let $S$ be a multiplicatively stable
subset of $A-(0)$ which contains $1$. Set $A'=S^{-1}A$.
\begin{enumerate}[label=(\arabic*),leftmargin=2.0em,itemsep=0.1ex,topsep=0.3ex]
\item\label{prop:V.1.1.1} For any ideal $\mathfrak{b}'$ of $A'$, it is true that
$(\mathfrak{b}'\cap A)A'=\mathfrak{b}'$, so the mapping
$\mathfrak{b}'\mapsto\mathfrak{b}'\cap A$ is an increasing (for the inclusion
relation) injection of the set of ideals of $A'$ into the set of ideals of
$A$.
\item\label{prop:V.1.1.2} The mapping $\mathfrak{p}'\mapsto\mathfrak{p}'\cap A$ is an isomorphism
of the partially ordered set (under inclusion) of prime ideals of $A'$ on the
partially ordered set of prime ideals $\mathfrak{p}$ of $A$ which satisfy
$\mathfrak{p}\cap S=\phi$. The inverse mapping is
$\mathfrak{p}\mapsto\mathfrak{p}A'$.
\end{enumerate}
\end{proposition}

\begin{proof}
1.\quad If $\mathfrak{b}'$ is an ideal of $A'$, then
$\mathfrak{b}'\cap A\subset\mathfrak{b}'$ and
$(\mathfrak{b}'\cap A)A'\subset\mathfrak{b}'$, since $\mathfrak{b}'$ is an
ideal. To prove the reverse inclusion take $x\in\mathfrak{b}'$. Since
$x=a/s$ ($a\in A$ and $s\in S$), $sx\in\mathfrak{b}'$ (since $A\subset A'$
and $\mathfrak{b}'$ is an ideal of $A'$). It follows that
$a\in\mathfrak{b}'$ and $a\in\mathfrak{b}'\cap A$. Therefore,
$x=\frac{1}{s}a\in A'(\mathfrak{b}'\cap A)$, so
$\mathfrak{b}'\subset A'(\mathfrak{b}'\cap A)$ and
$\mathfrak{b}'=A'(\mathfrak{b}'\cap A)$. This formula assures the injectivity
of the mapping $\varphi:\mathfrak{b}'\mapsto\mathfrak{b}'\cap A$, inasmuch as
there is a mapping $\theta:\mathfrak{b}\mapsto A'\mathfrak{b}$ such that
$\theta\circ\varphi=\text{identity}$. It is clear that $\varphi$ is an
increasing mapping.

2.\quad If $\mathfrak{p}'$ is a prime ideal of $A'$, then
$\mathfrak{p}=\mathfrak{p}'\cap A$ is a prime ideal of $A$ (\hyperref[sec:3.3]{Chapter III, §3},
\hyperref[lem:III.3.1]{Lemma 1}). Furthermore, $\mathfrak{p}\cap S=\phi$, since, if
$s\in\mathfrak{p}\cap S$, then $s\in\mathfrak{p}'$ and
$1=(1/s)s\in A'\mathfrak{p}'=\mathfrak{p}'$, an absurdity. Conversely, let
$\mathfrak{p}$ be a prime ideal of $A$ such that $\mathfrak{p}\cap S=\phi$.
We are going to show that $\mathfrak{p}A'$ is a prime ideal of $A'$ and that
$\mathfrak{p}A'\cap A=\mathfrak{p}$. Note first that $\mathfrak{p}A'$ is the
set of all elements of $A'$ which are of the form $p/s$ with
$p\in\mathfrak{p}$ and $s\in S$; any element $x\in\mathfrak{p}A'$ may be
written as
\[
x=\sum_{i=1}^{n}\frac{a_{i}}{s_{i}}p_{i}
\quad(a_{i}\in A,\ s_{i}\in S,\ \text{and } p_{i}\in\mathfrak{p}),
\]
so
\[
x=\sum_{i=1}^{n}\frac{b_{i}}{s}p_{i}
\left(s=s_{1}\ldots s_{n}\ \text{and}\ \frac{b_{i}}{s}=\frac{a_{i}}{s_{i}}\right).
\]
Thus $x=p/s$ with $p=\sum b_{i}p_{i}\in\mathfrak{p}$. We may conclude that
$1\notin\mathfrak{p}A'$, since $\mathfrak{p}\cap S=\phi$ and since $1$ cannot
be written in the form $p/s$ with $p\in\mathfrak{p}$ and $s\in S$. To show
that the ideal $\mathfrak{p}A'$ is prime let $a/s\in A'$ and $b/t\in A'$
with $(a/s)\cdot(b/t)\in\mathfrak{p}A'$. Then $(a/s)(b/t)=p/u$ with
$p\in\mathfrak{p}$ and $u\in S$, $abu=pst\in\mathfrak{p}$. Since
$\mathfrak{p}\cap S=\phi$, $u\notin\mathfrak{p}$, so
$ab\in\mathfrak{p}$ ($\mathfrak{p}$ is prime), which implies that either $a$
or $b$ belongs to $\mathfrak{p}$. Therefore, either $a/s$ or $b/t$ belongs to
$\mathfrak{p}A'$, i.e.\ $\mathfrak{p}A'$ is prime. Let us show finally that
$\mathfrak{p}=\mathfrak{p}A'\cap A$. Clearly,
$\mathfrak{p}\subset\mathfrak{p}A'\cap A$. For the reverse inclusion take
$x\in\mathfrak{p}A'\cap A$; then $x=p/s$ ($p\in\mathfrak{p}$, $s\in S$),
since, by hypothesis, $x\in\mathfrak{p}A'$. Thus $sx=p\in\mathfrak{p}$; since
$s\notin\mathfrak{p}$ ($\mathfrak{p}\cap S=\phi$) and since
$\mathfrak{p}$ is prime, it follows that $x\in\mathfrak{p}$. Now we simply
observe that the formulas $\mathfrak{p}=\mathfrak{p}A'\cap A$ and
$\mathfrak{p}'=A'(\mathfrak{p}'\cap A)$ entail that the mappings
$\varphi:\mathfrak{p}'\mapsto\mathfrak{p}'\cap A$ and
$\theta:\mathfrak{p}\mapsto\mathfrak{p}A'$ (restricted to the prime ideals
$\mathfrak{p}$ which do not intersect $S$) are inverse bijections, since
$\theta\circ\varphi$ and $\varphi\circ\theta$ are both identity mappings.
\end{proof}

\begin{statement}{Corollary}
If $A$ is a Noetherian integral domain, then every ring of fractions
$S^{-1}A$ is Noetherian.
\end{statement}

\begin{proof}
By Proposition \ref{prop:V.1.1}, \hyperref[prop:V.1.1.1]{(1)} there is an injective mapping of the
lattice of ideals in $S^{-1}A$ to a sublattice of the lattice of ideals in
$A$. Therefore, the lattice of ideals in $S^{-1}A$ satisfies the maximal
condition, so $S^{-1}A$ is a Noetherian ring.
\end{proof}

\begin{proposition}
\label{prop:V.1.2}
Let $R$ be an integral domain, let $A$ be a subring of $R$, let $S$ be a
multiplicatively stable subset of $A$ with $1\in S$ and $0\notin S$, and let
$B$ be the integral closure of $A$ in $R$. Then the integral closure of
$S^{-1}A$ in $S^{-1}R$ is $S^{-1}B$.
\end{proposition}

\begin{proof}
Any element of $S^{-1}B$ may be written in the form $b/s$ ($b\in B$,
$s\in S$). By dividing an equation of integral dependence for $b$ on $A$,
e.g.
\[
b^{n}+a_{n-1}b^{n-1}+\cdots+a_{0}=0\quad(a_{i}\in A),
\]
by $s^{n}$ we obtain an equation
\[
\left(\frac{b}{s}\right)^{n}+\frac{a_{n-1}}{s}
\left(\frac{b}{s}\right)^{n-1}+\cdots+\frac{a_{0}}{s^{n}}=0,
\]
which shows that $b/s$ is integral over $S^{-1}A$. Conversely, if $x/s$
($x\in R$, $s\in S$), an element of $S^{-1}R$, is integral over $S^{-1}A$,
then there is an equation of the form
\[
\left(\frac{x}{s}\right)^{n}+\frac{a_{n-1}}{t_{n-1}}
\left(\frac{x}{s}\right)^{n-1}+\cdots+\frac{a_{0}}{t_{0}}=0
\quad(a_{i}\in A,\ t_{i}\in S).
\]
Multiplying through by $(t_{0}t_{1}\ldots t_{n-1})^{n}$ shows that
$x t_{0}\ldots t_{n-1}/s$ is integral over $A$. Therefore,
$x t_{0}\ldots t_{n-1}/s\in B$ and
$x/s=(1/t_{0}\ldots t_{n-1})(x t_{0}\ldots t_{n-1}/s)$ belongs to
$S^{-1}B$.
\end{proof}

\begin{statement}{Corollary}
If $A$ is an integrally closed ring, then every ring of fractions $S^{-1}A$
is integrally closed.
\end{statement}

\begin{proof}
Take for $R$ in Proposition \ref{prop:V.1.2} the field of fractions of $A$.
\end{proof}

\begin{proposition}
\label{prop:V.1.3}
If $A$ is a Dedekind ring, every ring of fractions $S^{-1}A$ is a Dedekind
ring.
\end{proposition}

\begin{proof}
$S^{-1}A$ is Noetherian (the corollary of Proposition \ref{prop:V.1.1}) and
integrally closed (the corollary of Proposition \ref{prop:V.1.2}).
Furthermore, since one ``loses'' some prime ideals in passing from $A$ to
$S^{-1}A$ (Proposition \ref{prop:V.1.1}, \hyperref[prop:V.1.1.2]{(2)}), every non-zero prime ideal of
$S^{-1}A$ is maximal.
\end{proof}

\begin{proposition}
\label{prop:V.1.4}
Let $A$ be a Dedekind ring, let $\mathfrak{p}$ be a non-zero prime ideal of
$A$, and let $S=A-\mathfrak{p}$. Then $S^{-1}A$ is a principal ideal ring.
More precisely, there exists a prime $p\in S^{-1}A$ such that the only
non-zero ideals of $S^{-1}A$ are of the form $(p^{n})$, $n\geq0$.
\end{proposition}

\begin{proof}
Since $\mathfrak{p}$ is the only non-zero prime ideal of $A$ disjoint from
$S$, the only non-zero prime ideal of $S^{-1}A$ is
$\mathfrak{P}=\mathfrak{p}S^{-1}A$ (Proposition \ref{prop:V.1.1}, \hyperref[prop:V.1.1.2]{(2)}).
Since $S^{-1}A$ is a Dedekind ring (Proposition \ref{prop:V.1.3}), its only
non-zero ideals are of the form $\mathfrak{P}^{n}$ ($n\geq0$). Take an
element $p\in\mathfrak{P}-\mathfrak{P}^{2}$. The ideal $(p)$ is contained in
$\mathfrak{P}$ but not in $\mathfrak{P}^{2}$. Therefore,
$(p)=\mathfrak{P}$ and $(p^{n})=\mathfrak{P}^{n}$ for all $n\geq0$. Thus
$S^{-1}A$ is principal and all its ideals are of the form $(p^{n})$,
$n\geq0$.
\end{proof}

\begin{proposition}
\label{prop:V.1.5}
Let $A$ be an integral domain, $S$ a multiplicatively stable subset of
$A-(0)$ containing $1$, and let $\mathfrak{m}$ be a maximal ideal of $A$
which is disjoint from $S$. Then $S^{-1}A/\mathfrak{m}S^{-1}A\cong A/\mathfrak{m}$.
\end{proposition}

\begin{proof}
The composition of ring homomorphisms
$A\to S^{-1}A\to S^{-1}A/\mathfrak{m}S^{-1}A$ has kernel
$\mathfrak{m}S^{-1}A\cap A=\mathfrak{m}$, so there is an injection
$\varphi:A/\mathfrak{m}\to S^{-1}A/\mathfrak{m}S^{-1}A$. We must show that
$\varphi$ is surjective. Take $x=a/s\in S^{-1}A$ ($a\in A$, $s\in S$) and let
$\bar x$ denote its residue class in $S^{-1}A/\mathfrak{m}S^{-1}A$. Since
$s\notin\mathfrak{m}$ (by hypothesis, $\mathfrak{m}\cap S=\phi$) and since
$\mathfrak{m}$ is maximal, there exists $b\in A$ such that
$bs\equiv1\pmod{\mathfrak{m}}$. Thus
\[
\frac{a}{s}-ab=\frac{a}{s}(1-bs)\in\mathfrak{m}S^{-1}A,
\]
so $\varphi(ab)=\bar x$.
\end{proof}

%% ----------------------------------------------------------------------
\sechead{2}{The splitting of a prime ideal in an extension}%
        {THE SPLITTING OF A PRIME IDEAL IN AN EXTENSION}

In this section $A$ denotes a Dedekind ring of characteristic zero, $K$ its
field of fractions, $L$ a finite extension of $K$ of degree $n$, and $B$ the
integral closure of $A$ in $L$. We remind the reader that $B$ is also a
Dedekind ring (\hyperref[sec:3.4]{Chapter III, §4}, Theorem \ref{thm:III.4.1}).

Let $\mathfrak{p}$ be a non-zero prime ideal of $A$. Then
$B\mathfrak{p}$ is an ideal of $B$ and it has an expression of the form
\begin{equation}
B\mathfrak{p}=\prod_{i=1}^{q}\mathfrak{P}_{i}^{e_{i}},
\label{eq:V.2.1}
\end{equation}
where the $\mathfrak{P}$'s are distinct prime ideals of $B$, the $e_{i}$'s
are positive integers, and the product sign denotes multiplication of
ideals.

\begin{proposition}
\label{prop:V.2.1}
The $\mathfrak{P}_{i}$'s are precisely those prime ideals $\Sigma$ of $B$
such that $\Sigma\cap A=\mathfrak{p}$.
\end{proposition}

\begin{proof}
For a prime ideal $\Sigma$ of $B$ the relation $\Sigma\cap A=\mathfrak{p}$
is equivalent to the relation $\Sigma\supset B\mathfrak{p}$
($\Rightarrow$ is clear, $\Leftarrow$ follows from the fact that
$\Sigma\cap A$ is a prime ideal of $A$ and $\mathfrak{p}$ is maximal).
Clearly, $B\mathfrak{p}=\prod_{i=1}^{q}\mathfrak{P}_{i}^{e_{i}}$ implies
$B\mathfrak{p}\subset\mathfrak{P}_{i}$ for each $i=1,\ldots,q$, so
$\mathfrak{P}_{i}$ appears in the product expression for $B\mathfrak{p}$ if
and only if $\mathfrak{P}_{i}\cap A=\mathfrak{p}$.
\end{proof}

It is now clear that $A/\mathfrak{p}$ may be identified with a subring of
$B/\mathfrak{P}_{i}$ for any $i=1,\ldots,q$. Both rings are fields. Since $B$
is an $A$-module of finite type (\hyperref[sec:3.4]{Chapter III, §4}, Theorem
\ref{thm:III.4.1}), $B/\mathfrak{P}_{i}$ is a finite dimensional vector
space over $A/\mathfrak{p}$. We shall write $f_{i}$ for the dimension of
$B/\mathfrak{P}_{i}$ over $A/\mathfrak{p}$ and call $f_{i}$ the
\emph{residual degree} of $\mathfrak{P}_{i}$ over $A$. The exponent $e_{i}$
in \eqref{eq:V.2.1} is called the \emph{ramification index} of
$\mathfrak{P}_{i}$ over $A$. Let us remark finally that
$B\mathfrak{p}\cap A=\mathfrak{p}$ ($\supset$ clear, $\subset$ follows from
the fact that, for each $i=1,\ldots,q$,
$\mathfrak{P}_{i}\cap A=\mathfrak{p}$), so $B/B\mathfrak{p}$ is a vector
space over $A/\mathfrak{p}$, also of finite dimension.

\begin{theorem}
\label{thm:V.2.1}
With the preceding notations
\begin{equation}
\sum_{i=1}^{q}e_{i}f_{i}=[B/B\mathfrak{p}:A/\mathfrak{p}]=n.
\label{eq:V.2.2}
\end{equation}
\end{theorem}

\begin{proof}
The first equality is easy. Consider the sequence of ideals
\[
B\supset\mathfrak{P}_{1}\supset\mathfrak{P}_{1}^{2}\supset\cdots\supset
\mathfrak{P}_{1}^{e_{1}}\supset
\mathfrak{P}_{1}^{e_{1}}\mathfrak{P}_{2}\supset\cdots\supset
\mathfrak{P}_{1}^{e_{1}}\mathfrak{P}_{2}^{e_{2}}\supset\cdots\supset
\mathfrak{P}_{1}^{e_{1}}\cdots\mathfrak{P}_{q}^{e_{q}}=B\mathfrak{p}.
\]
Two consecutive elements of this sequence are of the form $\mathfrak{B}$ and
$\mathfrak{B}\mathfrak{P}_{i}$. Since there is no ideal strictly between
$\mathfrak{B}$ and $\mathfrak{B}\mathfrak{P}_{i}$,
$\mathfrak{B}/\mathfrak{B}\mathfrak{P}_{i}$ is a vector space of dimension
one over $B/\mathfrak{P}_{i}$ (cf.\ the proof of Proposition, \hyperref[sec:3.5]{§5, Chapter
III}). Thus it is a vector space of dimension $f_{i}$ over
$A/\mathfrak{p}$. For a given $i$ there are exactly $e_{i}$ consecutive
elements of the above sequence with associated quotient space of the form
$\mathfrak{B}/\mathfrak{B}\mathfrak{P}_{i}$, i.e.\ of dimension $f_{i}$ over
$A/\mathfrak{p}$. The total dimension
$[B/B\mathfrak{p}:A/\mathfrak{p}]$ equals the sum of the dimensions of the
quotients, so it is $\sum_{i=1}^{q}e_{i}f_{i}$.

The second equality is also easy in the case where $B$ is a free
$A$-module, in particular when $A$ is principal (\hyperref[sec:2.7]{Chapter II, §7}, the
corollary of Theorem \ref{thm:II.7.1}). In this case a base
$(x_{1},\ldots,x_{n})$ of $B$ as an $A$-module gives, by reduction mod
$B\mathfrak{p}$, a base for $B/B\mathfrak{p}$ over $A/\mathfrak{p}$. We are
going to reduce the general case to this case by considering the
multiplicatively stable subset $S=A-\mathfrak{p}$ of $A$ and the rings of
fractions $A'=S^{-1}A$ and $B'=S^{-1}B$. We know that $A'$ is a principal
ideal ring in which $\mathfrak{p}A'$ is the unique maximal ideal (\hyperref[sec:5.1]{§1},
Proposition \ref{prop:V.1.4}), and that $B'$ is the integral closure of $A'$
in $L$ (\hyperref[sec:5.1]{§1}, Proposition \ref{prop:V.1.2}). By the special case when $A$ is
principal, $[B'/\mathfrak{p}B':A'/\mathfrak{p}A']=n$. Now consider the
factorization of the ideal $\mathfrak{p}B'$ in the Dedekind ring $B'$: from
the fact that
$B\mathfrak{p}=\prod_{i=1}^{q}\mathfrak{P}_{i}^{e_{i}}$, we conclude that
$\mathfrak{p}B'=\prod_{i=1}^{q}(B'\mathfrak{P}_{i})^{e_{i}}$. Since
$\mathfrak{P}_{i}\cap A=\mathfrak{p}$ (Proposition \ref{prop:V.2.1}),
$\mathfrak{P}_{i}\cap S=\phi$ and $B'\mathfrak{P}_{i}$ is a non-zero prime
ideal of $B'$ (\hyperref[sec:5.1]{§1}, Proposition \ref{prop:V.1.1}, \hyperref[prop:V.1.1.2]{(2)}). From the first part
of our proof we now obtain
\[
[B'/\mathfrak{p}B':A'/\mathfrak{p}A']
=\sum_{i=1}^{q}e_{i}[B'/B'\mathfrak{P}_{i}:A'/\mathfrak{p}A'].
\]
However,
\[
A'/\mathfrak{p}A'\cong A/\mathfrak{p}\quad\text{and}\quad
B'/B'\mathfrak{P}_{i}\cong B/\mathfrak{P}_{i}\quad
(\S1,\ \text{Proposition \ref{prop:V.1.5}})
\]
Therefore,
\[
n=[B'/\mathfrak{p}B':A'/\mathfrak{p}A']=\sum_{i=1}^{q}e_{i}f_{i}.
\]
\end{proof}

\begin{proposition}
\label{prop:V.2.2}
With the same notations, the ring $B/B\mathfrak{p}$ is isomorphic to the
ring $\prod_{i=1}^{q}B/\mathfrak{P}_{i}^{e_{i}}$.
\end{proposition}

\begin{proof}
$\mathfrak{P}_{i}$ is the only maximal ideal of $B$ which contains
$\mathfrak{P}_{i}^{e_{i}}$, so
$\mathfrak{P}_{i}^{e_{i}}+\mathfrak{P}_{j}^{e_{j}}=B$ for $i\neq j$. The
proposition now follows from \eqref{eq:V.2.1} and \hyperref[lem:I.3.1]{Lemma 1}, \hyperref[sec:1.3]{§3, Chapter I}.
\end{proof}

\begin{statement}{Example}
Cyclotomic fields. Let $p$ be a prime number and let $z$ be a
primitive $p^{r}$th root of unity in the complex numbers. In this case, all
the complex $p^{r}$th roots of unity are of the form $z^{j}$
($j=1,\ldots,p^{r}$). The primitive roots of unity are those for which $j$
is not a multiple of $p$. The number of primitive roots is
$\varphi(p^{r})=p^{r}-p^{r-1}=p^{r-1}(p-1)$ (cf.\ \hyperref[sec:1.6]{Chapter I, §6}). The
primitive $p^{r}$th roots of unity are the roots of the cyclotomic
polynomial
\begin{equation}
F(X)=\frac{X^{p^{r}}-1}{X^{p^{r-1}}-1}
=X^{p^{r-1}(p-1)}+X^{p^{r-1}(p-2)}+\cdots+X^{p^{r-1}}+1.
\label{eq:V.2.3}
\end{equation}
We intend to give another proof that $[\Q[z]:\Q]=p^{r-1}(p-1)$, i.e.\ that
$F$ is irreducible (cf.\ \hyperref[sec:2.9]{Chapter II, §9}). For this purpose put
$e=p^{r-1}(p-1)$ and let $z_{1},\ldots,z_{e}$ be all the primitive
$p^{r}$th roots of unity. Since the constant term of $F(X+1)$ is $p$,
$\prod_{j=1}^{e}(z_{j}-1)=\pm p$. Let $B$ be the ring of integers of
$\Q[z]$. Clearly, $z_{j}\in B$ and
$z_{j}-1\in B(z_{k}-1)$, for all $j$ and $k$, since $z_{j}$ is a power
$z_{k}^{q}$ of $z_{k}$ and
$z_{k}^{q}-1=(z_{k}-1)(z_{k}^{q-1}+\cdots+z_{k}+1)$. Thus all the ideals
$B(z_{k}-1)$ are the same. It follows that $Bp=B(z_{1}-1)^{e}$.

Now write $Bp=\prod_{i=1}^{q}\mathfrak{P}_{i}^{e_{i}}$, where the
$\mathfrak{P}_{i}$'s are prime ideals of $B$. The $e_{i}$'s must, clearly,
all be multiples of $e$. But $e\geq[\Q[z]:\Q]$ (by \eqref{eq:V.2.3}), so
$e\geq\sum_{i=1}^{q}e_{i}f_{i}$ (Theorem \ref{thm:V.2.1}). From these
inequalities (they are really equalities) we may conclude that $q=1$,
$e=e_{1}$, $f_{1}=1$, and $[\Q[z]:\Q]=e$. In summary:
\begin{alist}
\item $[\Q[z]:\Q]=e=p^{r-1}(p-1)$,
\item $B(z_{1}-1)$ is a prime ideal of $B$ of residual degree $1$, and
\item $Bp=B(z_{1}-1)^{e}$.
\end{alist}
\end{statement}

%% ----------------------------------------------------------------------
\sechead{3}{The discriminant and ramification}{THE DISCRIMINANT AND RAMIFICATION}

With the notations of \hyperref[sec:5.2]{§2} (let
$B\mathfrak{p}=\prod_{i=1}^{q}\mathfrak{P}_{i}^{e_{i}}$) a prime ideal
$\mathfrak{p}$ of $A$ is said to \emph{ramify} in $B$ (or in $L$) if any one
of its ramification indices $e_{i}$ is larger than one. In terms of the
theory of the discriminant (\hyperref[sec:2.7]{Chapter II, §7}) we are going to characterize
those prime ideals of $A$ which ramify in $B$. In particular we shall show
that only finitely many prime ideals of $A$ ramify in $B$. First we need
some lemmas concerning the discriminant.

\begin{lemma}
\label{lem:V.3.1}
Let $A$ be a ring, let $B_{1},\ldots,B_{q}$ be rings containing $A$ which
are free $A$-modules of finite type, and let
$B=\prod_{i=1}^{q}B_{i}$ be the product ring. Then
$\mathfrak{D}_{B/A}=\prod_{i=1}^{q}\mathfrak{D}_{B_{i}/A}$ (cf.\ \hyperref[sec:2.7]{Chapter II,
§7}, Definition \ref{def:II.7.2}).
\end{lemma}

\begin{proof}
We formulate our proof for the case $q=2$. The general statement follows by
induction on $q$. Let $(x_{1},\ldots,x_{m})$ and $(y_{1},\ldots,y_{n})$ be
bases for $B_{1}$ and $B_{2}$ as modules over $A$. With the usual
identifications of $B_{1}$ with $B_{1}\times(0)$ and $B_{2}$ with
$(0)\times B_{2}$, we may consider $(x_{1},\ldots,x_{m},y_{1},\ldots,y_{n})$
as a base for $B=B_{1}\times B_{2}$ over $A$. By definition of the product
ring structure $x_{i}y_{j}=0$, from which it follows that
$\operatorname{Tr}(x_{i}y_{j})=0$. As a consequence the determinant
$\mathrm{D}(x_{1},\ldots,x_{m},y_{1},\ldots,y_{n})$ is the determinant of
the matrix
\[
\begin{vmatrix}
\operatorname{Tr}(x_{i}x_{i'}) & 0\\[2pt]
0 & \operatorname{Tr}(y_{j}y_{j'})
\end{vmatrix}
\]
The value of this determinant is
\[
\det(\operatorname{Tr}(x_{i}x_{i'}))
\det(\operatorname{Tr}(y_{j}y_{j'})).
\]
so
\[
\mathrm{D}(x_{1},\ldots,x_{m},y_{1},\ldots,y_{n})
=\mathrm{D}(x_{1},\ldots,x_{m})\mathrm{D}(y_{1},\ldots,y_{r}).
\]
\end{proof}

\begin{lemma}
\label{lem:V.3.2}
Let $B$ be a ring, $A$ a subring of $B$, and $\mathfrak{a}$ an ideal of $A$.
Assume that $B$ is a free module over $A$ with the base
$(x_{1},\ldots,x_{n})$. For $x\in B$ write $\bar x$ for the residue class of
$x$ in $B/\mathfrak{a}B$. Then $(\bar x_{1},\ldots,\bar x_{n})$ is a base of
$B/\mathfrak{a}B$ over $A/\mathfrak{a}$ and
\[
\mathrm{D}(\bar x_{1},\ldots,\bar x_{n})
=\overline{\mathrm{D}(x_{1},\ldots,x_{n})}.
\]
\end{lemma}

\begin{proof}
Let $x\in B$. If the matrix for multiplication by $x$, with respect to the
base $(x_{1},\ldots,x_{n})$ is $(a_{ij})$ ($a_{ij}\in A$), then the matrix
for multiplication by $\bar x$ with respect to the base
$(\bar x_{1},\ldots,\bar x_{n})$ is $(\bar a_{ij})$. Thus,
$\operatorname{Tr}(\bar x)=\overline{\operatorname{Tr}(x)}$. Taking
$x=x_{i}x_{j}$, we obtain
$\operatorname{Tr}(\bar x_{i}\bar x_{j})
=\overline{\operatorname{Tr}(x_{i}x_{j})}$, and the assertion follows by taking determinants.
\end{proof}

\begin{lemma}
\label{lem:V.3.3}
Let $K$ be a field which is finite or of characteristic zero. Let $L$ be a
finite dimensional (commutative) $K$-algebra. In order that $L$ be reduced
it is necessary and sufficient that $\mathfrak{D}_{L/K}\neq(0)$.
\end{lemma}

\begin{proof}
Suppose first that $L$ is not reduced and let $x$ be a non-zero nilpotent
element of $L$. Let $(x_{1},\ldots,x_{n})$ be a base for $L$ over $K$ such
that $x=x_{1}$. Then $x_{1}x_{j}$ is nilpotent and multiplication by
$x_{1}x_{j}$ is a nilpotent endomorphism of the vector space $L$ over $K$.
Thus, all the characteristic values of this endomorphism are zero, so
$\operatorname{Tr}(x_{1}x_{j})=0$.

The matrix $(\operatorname{Tr}(x_{i}x_{j}))$ has a row comprised entirely of
zeroes, which implies that $\mathrm{D}(x_{1},\ldots,x_{n})=0$, i.e.\
$\mathfrak{D}_{L/K}=(0)$.

Conversely, suppose that $L$ is reduced. Then the ideal $(0)$ of $L$ is
expressible as a finite intersection of prime ideals,
$(0)=\bigcap_{i=1}^{q}\mathfrak{P}_{i}$ (\hyperref[sec:4.7]{Chapter IV, §7}, Lemma). Since
$L/\mathfrak{P}_{i}$ is an integral domain and a finite dimensional algebra
over $K$, it is a field (\hyperref[sec:2.1]{Chapter II, §1}, Proposition \ref{prop:II.1.3}). It
follows that $\mathfrak{P}_{i}$ is a maximal ideal of $L$, and consequently
$\mathfrak{P}_{i}+\mathfrak{P}_{j}=L$ for $i\neq j$. Therefore, $L$ is
isomorphic to the ring product $\prod_{i=1}^{q}L/\mathfrak{P}_{i}$ (\hyperref[sec:1.3]{Chapter
I, §3}, Lemma \ref{lem:I.3.1}). By Lemma \ref{lem:V.3.1}
$\mathfrak{D}_{L/K}=\prod_{i=1}^{q}\mathfrak{D}_{(L/\mathfrak{P}_{i})/K}$.
But $\mathfrak{D}_{(L/\mathfrak{P}_{i})/K}\neq(0)$, since $K$ is finite or
of characteristic zero (\hyperref[sec:2.7]{Chapter II, §7}, Proposition \ref{prop:II.7.3}).
Therefore, $\mathfrak{D}_{L/K}\neq(0)$.
\end{proof}

\begin{definition}
\label{def:V.3.1}
Let $K$ and $L$ be number fields with $K\subset L$. Let $A$ and $B$ be the
rings of integers of $K$ and $L$, respectively. The \emph{discriminant}
(ideal) of $B$ over $A$ (or of $L$ over $K$) is the ideal of $A$ generated
by the discriminants of bases of $L$ over $K$ which are contained in $B$.
Notation: $\mathfrak{D}_{B/A}$ or $\mathfrak{D}_{L/K}$.
\end{definition}

\begin{remark}
\label{rem:V.3.1}
If $(x_{1},\ldots,x_{n})$ is a base of $L$ over $K$ contained in $B$, then
$\operatorname{Tr}_{L/K}(x_{i}x_{j})\in A$ (\hyperref[sec:2.6]{Chapter II, §6}, the corollary of
Proposition \ref{prop:II.6.2}), so $\mathrm{D}(x_{1},\ldots,x_{n})\in A$.
Thus $\mathfrak{D}_{B/A}$ is an integral ideal of $A$. It is non-zero by
\hyperref[sec:2.7]{Chapter II, §7}, Proposition \ref{prop:II.7.3}.
\end{remark}

\begin{remark}
\label{rem:V.3.2}
When $B$ is a free $A$-module (for example if $A$ is principal) we have
already defined the discriminant $\mathfrak{D}_{B/A}$ as the ideal generated
by $\mathrm{D}(e_{1},\ldots,e_{n})$ where $(e_{1},\ldots,e_{n})$ is an
$A$-module base for $B$ (\hyperref[sec:2.7]{Chapter II, §7}, Definition \ref{def:II.7.2}). Our
old definition coincides with the one given above, since, given any base
$(x_{i})$ of $L$ over $K$ contained in $B$, one sees that
$x_{i}=\sum_{j}a_{ij}e_{j}$ with $a_{ij}\in A$. Therefore,
\[
\mathrm{D}(x_{1},\ldots,x_{n})
=\det(a_{ij})^{2}\mathrm{D}(e_{1},\ldots,e_{n})
\]
(\hyperref[sec:2.7]{Chapter II, §7}, Proposition \ref{prop:II.7.1}).
\end{remark}

\begin{theorem}
\label{thm:V.3.1}
Let the notations be as in the definition. In order that a prime ideal
$\mathfrak{p}$ of $A$ ramify in $B$, it is necessary and sufficient that it
contain the discriminant $\mathfrak{D}_{B/A}$. There are only finitely many
prime ideals of $A$ which ramify in $B$.
\end{theorem}

\begin{proof}
The second assertion follows from the first, since
$\mathfrak{D}_{B/A}\neq(0)$. Let us prove the first. Since
$B/\mathfrak{p}B\cong\prod_{i=1}^{q}B/\mathfrak{P}_{i}^{e_{i}}$
(\hyperref[sec:5.2]{§2}, Proposition \ref{prop:V.2.2}), the assertion ``$\mathfrak{p}$
ramifies'' is equivalent to the statement ``$B/\mathfrak{p}B$ is not
reduced'', i.e.\ equivalent to
``$\mathfrak{D}_{(B/\mathfrak{p}B)/(A/\mathfrak{p})}=(0)$'' (by Lemma
\ref{lem:V.3.3} and the fact that $A/\mathfrak{p}$ is a finite field). Now
put $S=A-\mathfrak{p}$, $A'=S^{-1}A$, $B'=S^{-1}B$, and
$\mathfrak{p}'=\mathfrak{p}A'$. Then $A'$ is a principal ideal ring (\hyperref[sec:5.1]{§1},
Proposition \ref{prop:V.1.4}), $B'$ is a free $A'$-module,
$A/\mathfrak{p}\cong A'/\mathfrak{p}'$, and
$B/\mathfrak{p}B\cong B'/\mathfrak{p}'B'$ (\hyperref[sec:5.1]{§1}, Proposition
\ref{prop:V.1.5}). Therefore, writing $(e_{1},\ldots,e_{n})$ for an
$A'$-module base of $B'$, we know that
$\mathfrak{D}_{(B/\mathfrak{p}B)/(A/\mathfrak{p})}=(0)$ if and only if
$\mathrm{D}(e_{1},\ldots,e_{n})\in\mathfrak{p}'$ (Lemma \ref{lem:V.3.2}).
If $\mathrm{D}(e_{1},\ldots,e_{n})\in\mathfrak{p}'$ and if
$(x_{1},\ldots,x_{n})$ is a base for $L$ over $K$ contained in $B$, then
$x_{i}=\sum a'_{ij}e_{j}$ with $a'_{ij}\in A'$ (for $B\subset B'$), so
\[
\mathrm{D}(x_{1},\ldots,x_{n})
=\det(a'_{ij})^{2}\mathrm{D}(e_{1},\ldots,e_{n})\in\mathfrak{p}'.
\]
Since $\mathfrak{p}'\cap A=\mathfrak{p}$ (\hyperref[sec:5.1]{§1}, Proposition
\ref{prop:V.1.1}, \hyperref[prop:V.1.1.2]{(2)}), we may conclude that
$\mathrm{D}(x_{1},\ldots,x_{n})\in\mathfrak{p}$ and
$\mathfrak{D}_{B/A}\subset\mathfrak{p}$. Conversely, if
$\mathfrak{D}_{B/A}\subset\mathfrak{p}$, then
$\mathrm{D}(e_{1},\ldots,e_{n})\in\mathfrak{p}'$, for one may write
$e_{i}=y_{i}s^{-1}$ with $y_{i}\in B$ and $s\in S$, for $1\leq i\leq n$.
Consequently,
\[
\mathrm{D}(e_{1},\ldots,e_{n})=s^{-2n}\mathrm{D}(y_{1},\ldots,y_{n})
\in A'\mathfrak{D}_{B/A}\subset A'\mathfrak{p}=\mathfrak{p}'.
\]
\end{proof}

\begin{example}\label{exa:V.3.1}
Quadratic fields. Let $K=\Q$ and $L=\Q[\sqrt{d}]$, where $d$ is a
square-free integer (\hyperref[sec:2.5]{Chapter II, §5}).
\begin{alist}
\item If $d\equiv2$ or $3\pmod4$, then $(1,\sqrt{d})$ is a base for the
ring of integers of $L$. As $\operatorname{Tr}(1)=2$,
$\operatorname{Tr}(\sqrt{d})=0$, and $\operatorname{Tr}(d)=2d$, it follows
that $\mathrm{D}(1,\sqrt{d})=4d$. The prime numbers which ramify in $L$
therefore include $2$ and the prime divisors of $d$.
\item If $d\equiv1\pmod4$, $(1,(1+\sqrt{d})/2)$ is a base for the ring of
integers of $L$. We see that
\[
\operatorname{Tr}(1)=2,\quad
\operatorname{Tr}\left(\frac{1+\sqrt{d}}{2}\right)=1,
\]
and
\[
\operatorname{Tr}\left(\left(\frac{1+\sqrt{d}}{2}\right)^{2}\right)
=\operatorname{Tr}\left(\frac{d+1}{4}+\frac{1}{2}\sqrt{d}\right)
=\frac{d+1}{2},
\]
thus
\[
\mathrm{D}\left(1,\frac{1+\sqrt{d}}{2}\right)
=2\cdot\frac{d+1}{2}-1=d.
\]
The only prime numbers which ramify in $L$ are the divisors of $d$.
\end{alist}
\end{example}

We remark that a quadratic field $\Q[\sqrt{d}]$ is uniquely determined by its
discriminant. In fact,
\[
\mathrm{D}\equiv0\pmod4 \text{ implies } d=\frac{\mathrm{D}}{4}
\quad(\text{we must have } d\equiv2 \text{ or } 3\pmod4),
\]
\[
\mathrm{D}\equiv1\pmod4 \text{ implies } d=\mathrm{D},
\]
and $\mathrm{D}\equiv2$ or $3\pmod4$ is impossible.

\begin{example}\label{exa:V.3.2}
Cyclotomic fields. Let $p$ be an odd prime number, let $z$ be a
primitive complex $p$th root of unity, and let $L=\Q[z]$ be the
corresponding cyclotomic field. We know that $(1,z,\ldots,z^{p-2})$ is a
$\Z$-base for the ring of integers $B$ of $L$ (\hyperref[sec:2.9]{Chapter II, §9}, Theorem
\ref{thm:II.9.2}), and that the minimal polynomial $F(X)$ of $z$ over $\Q$
satisfies the relation $(X-1)F(X)=X^{p}-1$ (ibid., Theorem
\ref{thm:II.9.1}). We are going to calculate the discriminant
$\mathfrak{D}_{B/\Z}$ by making use of the formula
\[
\mathrm{D}(1,z,\ldots,z^{p-2})=(-1)^{\frac{1}{2}(p-1)(p-2)}
\mathrm{N}(F'(z))\quad(\text{\hyperref[sec:2.7]{Chapter II, §7}, formula
\eqref{eq:II.7.6}}).
\]
By taking the derivative of both sides of the relation $(X-1)F(X)=X^{p}-1$,
we obtain $(z-1)F'(z)=pz^{p-1}$ (since $F(z)=0$). We know that
$\mathrm{N}(p)=p^{p-1}$, $\mathrm{N}(z)=\pm1$, and
$\mathrm{N}(z-1)=\pm p$ (\hyperref[sec:2.9]{Chapter II, §9}). Therefore,
\begin{equation}
\mathrm{D}(1,z,\ldots,z^{p-2})=\pm p^{p-2}.
\label{eq:V.3.1}
\end{equation}
It follows that $p$ is the only prime number which ramifies in $\Q[z]$.
\end{example}

The following result is sometimes useful for determining the ring of
integers of a number field.

\begin{proposition}
\label{prop:V.3.1}
Let $L$ be a number field of degree $n$ over $\Q$ and let
$(x_{1},\ldots,x_{n})$ be a $\Q$-base for $L$ contained in the ring $B$ of
integers of $L$. If the discriminant $\mathrm{D}(x_{1},\ldots,x_{n})$ is
square-free, then $(x_{1},\ldots,x_{n})$ is a $\Z$-base for $B$.
\end{proposition}

\begin{proof}
If $(e_{1},\ldots,e_{n})$ is a $\Z$-base for $B$, then
$x_{i}=\sum_{j=1}^{n}a_{ij}e_{j}$, with $a_{ij}\in\Z$, whence
\[
\mathrm{D}(x_{1},\ldots,x_{n})
=\det(a_{ij})^{2}\mathrm{D}(e_{1},\ldots,e_{n}).
\]
Since $\mathrm{D}(x_{1},\ldots,x_{n})$ is square-free, $\det(a_{ij})$ must
be $\pm1$, so $(x_{1},\ldots,x_{n})$ is also a $\Z$-base for $B$.
\end{proof}

\begin{statement}{Remark}
The cyclotomic fields (for $p\geq5$), or quadratic fields, provide examples
which show that the sufficient condition of Proposition \ref{prop:V.3.1} is
not a necessary condition.
\end{statement}

\begin{statement}{Example}
The polynomial $X^{3}-X-1$ (respectively, $X^{3}+X+1$, $X^{3}+10X+1$) is
irreducible over $\Q$. Otherwise it would have a linear factor, i.e.\ a
rational root $x\in\Q$. But, the polynomial is monic, so $x\in\Q$ implies
$x\in\Z$. The constant term of the polynomial is $1$, so $x=\pm1$ ($x$ has
to divide the constant term). Checking that neither $\pm1$ satisfies the
polynomial, we may conclude that the polynomial is irreducible.

Now let $x$ be a complex root of the given polynomial. The field
$L=\Q[x]$ is a cubic field (i.e.\ of degree $3$). Thus, $(1,x,x^{2})$ is a
base for $L$ over $\Q$, and $x$ is, clearly, an integer of $L$. By formula
\eqref{eq:II.7.7} of \hyperref[sec:2.7]{§7, Chapter II},
$\mathrm{D}(1,x,x^{2})=4-27=-23$ (respectively $-31$, $-4027$), the
negative of a prime number. Therefore, $(1,x,x^{2})$ is a $\Z$-base for the
ring of integers of $L$.
\end{statement}

%% ----------------------------------------------------------------------
\sechead{4}{The splitting of a prime number in a quadratic field}%
        {THE SPLITTING OF A PRIME NUMBER IN A QUADRATIC FIELD}

Let $d\in\Z$ be square-free, let $L$ be the quadratic field $\Q[\sqrt{d}]$,
write $B$ for the ring of integers of $L$, and let $p$ be a prime number. We
are going to study the factorization of the ideal $pB$ into a product of
prime ideals of $B$.

The formula $\sum_{i=1}^{q}e_{i}f_{i}=2$ (\hyperref[sec:5.2]{§2}, Theorem \ref{thm:V.2.1})
entails $q\leq2$ and the following three possibilities:
\begin{alist}
\item $q=2$, $e_{1}=e_{2}=1$, $f_{1}=f_{2}=1$;
in this case we say that $p$ \emph{splits} in $L$.
\item $q=1$, $e_{1}=1$, $f_{1}=2$;
in this case we say that $p$ \emph{remains prime} in $L$.
\item $q=1$, $e_{1}=2$, $f_{1}=1$;
this means that $p$ \emph{ramifies} in $L$.
\end{alist}

Let us first consider the case in which $p$ is odd. We know (\hyperref[sec:2.5]{Chapter II, §5})
that $B=\Z+\Z\sqrt{d}$ or $B=\Z+\Z[(1+\sqrt{d})/2]$, depending upon $d$.
But, if we pass to the residue classes of $B$ modulo $Bp$, we see in the
second case that $a+b[(1+\sqrt{d})/2]$ (with $b$ odd) is congruent to
$a+(b+p)[(1+\sqrt{d})/2]$, which belongs to $\Z+\Z\sqrt{d}$. Thus, for any
$d$, we have $B/Bp\cong(\Z+\Z\sqrt{d})/(p)$. Also we see that
\[
\Z+\Z\sqrt{d}\cong\Z[X]/(X^{2}-d),
\]
so
\[
B/Bp\cong\Z[X]/(p,X^{2}-d)\cong(\Z[X]/(p))/(X^{2}-d)
\cong\F_{p}[X]/(X^{2}-d),
\]
where $d$ denotes the residue class of $d$ modulo $p$. Now the assertion
that $p$ splits (respectively, remains prime, ramifies) in $B$ has the
interpretation: $B/Bp$ is the product of two fields (respectively, is a
field, contains nilpotent elements) (cf.\ \hyperref[sec:5.2]{§2}, Proposition
\ref{prop:V.2.2}). In other words, the polynomial $X^{2}-d\in\F_{p}[X]$ is
the product of two distinct linear polynomials (respectively, is
irreducible, is a square). This happens if $d$ is a non-zero square in
$\F_{p}$ (respectively, is not a square in $\F_{p}$, is zero in $\F_{p}$).
When $d$ is a non-zero square in $\F_{p}$ (respectively, is not a square in
$\F_{p}$), we say that $d$ is a \emph{quadratic residue} (respectively,
\emph{non-residue}) modulo $p$.

Let us now consider the case $p=2$. If $d\equiv2$ or $3\pmod4$, then
$B=\Z+\Z\sqrt{d}$, so, as above, $B/2B\cong\F_{2}[X]/(X^{2}-d)$. In this
case $X^{2}-d$ equals $X^{2}$ or $X^{2}-1=(X+1)^{2}$, in either case a
square. Thus $2$ ramifies in $B$. If $d\equiv1\pmod4$, $(1+\sqrt{d})/2$ has
$X^{2}-X-(d-1)/4$ as its minimal polynomial, so, as above,
$B/2B\cong\F_{2}[X]/(X^{2}-X-\delta)$, where $\delta$ is the residue class
mod $2$ of $(d-1)/4$. For $d\equiv1\pmod8$, $\delta=0$ and
$X^{2}-X-\delta=X(X-1)$, so that $2$ splits. For $d\equiv5\pmod8$,
$\delta=1$ and $X^{2}-X-\delta=X^{2}+X+1$, which is irreducible in
$\F_{2}[X]$, so $2$ remains prime. In summary, we have proved the following:

\begin{proposition}
\label{prop:V.4.1}
Let $L=\Q[\sqrt{d}]$, the quadratic field associated with the square-free
integer $d$.
\begin{alist}
\item The odd primes $p$ for which $d$ is a quadratic residue mod $p$ split
in $L$. So does $2$, if $d\equiv1\pmod8$.
\item The odd primes $p$ for which $d$ is not a quadratic residue mod $p$
remain prime in $L$. So does $2$, if $d\equiv5\pmod8$.
\item The odd prime divisors of $d$ ramify in $L$. So does $2$, if
$d\equiv2$ or $3\pmod4$.
\end{alist}
\end{proposition}

Part (c) was proved earlier, as an example in \hyperref[sec:5.3]{§3}.

%% ----------------------------------------------------------------------
\sechead{5}{The quadratic reciprocity law}{THE QUADRATIC RECIPROCITY LAW}

Given an \emph{odd} prime $p$ and an integer $d$ relatively prime to $p$ we
introduced in \hyperref[sec:5.4]{§4} the locution ``$d$ is a quadratic residue mod $p$''
(respectively, ``$d$ is a non-residue mod $p$'') as meaning that the residue
class of $d$ mod $p$ is a square (respectively, not a square) in $\F_{p}^{*}$.
Now we define the \emph{Legendre symbol} as follows:
\begin{equation}
\begin{cases}
\left(\dfrac{d}{p}\right)=+1, &
\text{if } d \text{ is a quadratic residue mod } p,\\[6pt]
\left(\dfrac{d}{p}\right)=-1, &
\text{if } d \text{ is a non-residue mod } p.
\end{cases}
\label{eq:V.5.1}
\end{equation}
It is understood that $\left(\frac{d}{p}\right)$ is defined only for integers
$d$ which are relatively prime to $p$, i.e.\ $d\in\Z-p\Z$. The
multiplicative group $\F_{p}^{*}$ being cyclic of even order $p-1$
(\hyperref[sec:1.7]{Chapter I, §7}, Theorem \ref{thm:I.7.1}), the squares in $\F_{p}^{*}$ form a
subgroup $(\F_{p}^{*})^{2}$ of index $2$, and
$\F_{p}^{*}/(\F_{p}^{*})^{2}$ is isomorphic to $\{+1,-1\}$. Clearly, the
Legendre symbol stands for the composition of the following homomorphisms:
\[
\Z-p\Z\to\F_{p}^{*}\to\F_{p}^{*}/(\F_{p}^{*})^{2}\simeq\{+1,-1\}.
\]
As a consequence there is the formula:
\begin{equation}
\left(\frac{ab}{p}\right)=\left(\frac{a}{p}\right)\left(\frac{b}{p}\right),
\quad a,b\in\Z-p\Z.
\label{eq:V.5.2}
\end{equation}

\begin{proposition}
\label{prop:V.5.1}
(``Euler's criterion'') If $p$ is an odd prime and if $a\in\Z-p\Z$, then
\[
\left(\frac{a}{p}\right)\equiv a^{(p-1)/2}\pmod p.
\]
\end{proposition}

\begin{proof}
Write $w$ for a primitive root mod $p$ (\hyperref[sec:1.7]{Chapter I, §7}). Then
$a\equiv w^{j}\pmod p$, with $0\leq j\leq p-2$, since the residue class
$\bar w$ of $w$ generates $\F_{p}^{*}$. Clearly, $a$ is a quadratic residue
if and only if $j$ is even. Therefore,
$\left(\frac{a}{p}\right)=(-1)^{j}$. On the other hand, $\F_{p}^{*}$
contains only one element of order $2$: this element can be written either
as $\bar u^{(p-1)/2}$ or $-1$. In $\Z$, we have
$-1\equiv w^{(p-1)/2}\pmod p$. Thus,
\[
\left(\frac{a}{p}\right)=(-1)^{j}\equiv w^{j(p-1)/2}
\equiv a^{(p-1)/2}\pmod p.
\]
\end{proof}

Now we are going to prove a famous theorem, which provides us with a
relation between the Legendre symbols for distinct odd primes.

\begin{theorem}
\label{thm:V.5.1}
(``The quadratic reciprocity law of Legendre-Gauss'') If $p$ and $q$ are
distinct odd prime numbers, then
\[
\left(\frac{p}{q}\right)\left(\frac{q}{p}\right)
=(-1)^{(p-1)(q-1)/4}.
\]
\end{theorem}

\begin{proof}
Let us consider, in a suitable extension of $\F_{q}$, a primitive $p$th root
of unity $w$. Since $w^{p}=1$, the expression $w^{x}$ is well defined for
$x\in\F_{p}$. We shall also consider the Legendre symbol
$\left(\frac{x}{p}\right)$ as defined on $x\in\F_{p}^{*}$, for
$\left(\frac{d}{p}\right)$ clearly depends only on the residue class of $d$
mod $p$. For $a\in\F_{p}^{*}$ consider the ``Gaussian sum''
\begin{equation}
\tau(a)=\sum_{x\in\F_{p}^{*}}\left(\frac{x}{p}\right)w^{ax}.
\label{eq:V.5.3}
\end{equation}
It is an element of an extension field of $\F_{q}$. Putting $ax=y$, we have
\[
\tau(a)=\sum_{y\in\F_{p}^{*}}\left(\frac{ya^{-1}}{p}\right)w^{y}
=\left(\frac{a^{-1}}{p}\right)\sum_{y\in\F_{p}^{*}}\left(\frac{y}{p}\right)
w^{y}\quad(\text{by \eqref{eq:V.5.2}}),
\]
so
\begin{equation}
\tau(a)=\left(\frac{a}{p}\right)\tau(1).
\label{eq:V.5.4}
\end{equation}

On the other hand, since we are working in a field of characteristic $q$ and
since $\left(\frac{x}{p}\right)\in\F_{q}$, we have
$\tau(1)^{q}=\sum_{x\in\F_{p}^{*}}\left(\frac{x}{p}\right)^{q}w^{qx}$, so,
identifying $q$ with its residue class mod $p$, we obtain
\begin{equation}
\tau(1)^{q}=\tau(q).
\label{eq:V.5.5}
\end{equation}
Let us calculate $\tau(1)^{2}$. We have
\[
\tau(1)^{2}=\sum_{x,y\in\F_{p}^{*}}
\left(\frac{x}{p}\right)\left(\frac{y}{p}\right)w^{x+y}.
\]
Putting $y=tx$, we see that
\[
\tau(1)^{2}
=\sum_{x,t\in\F_{p}^{*}}\left(\frac{x}{p}\right)^{2}\left(\frac{t}{p}\right)
w^{x(1+t)}
=\sum_{x,t}\left(\frac{t}{p}\right)w^{x(1+t)}
=\sum_{t\in\F_{p}^{*}}\left[\left(\frac{t}{p}\right)
\sum_{x\in\F_{p}^{*}}w^{x(1+t)}\right].
\]
If $w^{1+t}\neq1$, then $\sum_{f=0}^{p-1}(w^{1+t})^{f}=0$ by the formula for
summing a geometric series. Since $(w^{1+t})^{0}=1$, we have
\[
\sum_{x\in\F_{p}^{*}}w^{x(1+t)}=-1,
\quad\text{if } w^{1+t}\neq1.
\]
If $w^{1+t}=1$, then $\sum_{x\in\F_{p}^{*}}w^{x(1+t)}=p-1$; and this happens
if and only if $t=-1$, since $w$ is a primitive $p$th root of unity.
Therefore,
\[
\tau(1)^{2}=\left(\frac{-1}{p}\right)(p-1)
-\sum_{\substack{t\in\F_{p}^{*}\\t\neq-1}}\left(\frac{t}{p}\right).
\]
As there are as many squares as non-squares in $\F_{p}^{*}$,
\[
\sum_{t\in\F_{p}^{*}}\left(\frac{t}{p}\right)=0,\quad\text{so}\quad
\tau(1)^{2}=\left(\frac{-1}{p}\right)(p-1)+\left(\frac{-1}{p}\right)
=\left(\frac{-1}{p}\right)p.
\]
Using Euler's criterion (Proposition \ref{prop:V.5.1}), we obtain
\begin{equation}
\tau(1)^{2}=(-1)^{(p-1)/2}p.
\label{eq:V.5.6}
\end{equation}
Finally, using \eqref{eq:V.5.4} and \eqref{eq:V.5.5}, we see that
\[
\tau(1)^{q}=\tau(q)=\left(\frac{q}{p}\right)\tau(1).
\]
By \eqref{eq:V.5.6}, $\tau(1)\neq0$, so we may simplify:
\[
\tau(1)^{q-1}=\left(\frac{q}{p}\right).
\]
Use of \eqref{eq:V.5.6} again yields
\[
\left(\frac{q}{p}\right)=(\tau(1)^{2})^{(q-1)/2}
=(-1)^{(p-1)(q-1)/4}p^{(q-1)/2}.
\]
Since $p^{(q-1)/2}=\left(\frac{p}{q}\right)$ (Proposition
\ref{prop:V.5.1}) and since
$\left(\frac{q}{p}\right)=\left(\frac{p}{q}\right)^{-1}$, we have proved the
quadratic reciprocity law.
\end{proof}

\begin{proposition}
\label{prop:V.5.2}
(``The complementary formulas'') If $p$ is an odd prime, then
\[
\begin{array}{ll}
\text{(a)} & \left(\dfrac{-1}{p}\right)=(-1)^{(p-1)/2}\quad\text{and}\\[8pt]
\text{(b)} & \left(\dfrac{2}{p}\right)=(-1)^{(p^{2}-1)/8}.
\end{array}
\]
\end{proposition}

\begin{proof}
Relation (a) is a special case of Euler's criterion (Proposition
\ref{prop:V.5.1}). We need only prove (b). Note first that, since the
squares of $1,3,5$, and $7$ mod $8$ are $1,1,1$, and $1$, $p^{2}\equiv1
\pmod8$, so formula (b) makes sense. Let us remark next that, in the group
$H=\{1,3,5,7\}$ of units in the ring $\Z/8\Z$, $\{1,7\}$ is a subgroup $H'$
of index $2$. Put $\theta(x)=1$ for $x\in H'$ and $\theta(x)=-1$ for
$x\in H-H'$, so that $\theta(xy)=\theta(x)\theta(y)$ for $x,y\in H$. Let $w$
be a primitive eighth root of unity in an extension of $\F_{p}$. As in
Theorem \ref{thm:V.5.1} consider, for $a\in H$, the ``Gaussian sum''
\begin{equation}
\tau(a)=\sum_{x\in H}\theta(x)w^{ax}.
\label{eq:V.5.7}
\end{equation}
As in Theorem \ref{thm:V.5.1} we see that $\tau(a)=\theta(a)\tau(1)$ and
$\tau(1)^{p}=\tau(p)$ (identifying $p$ with its residue class mod $8$). From
the definition of $\theta(x)$ we obtain
\[
\tau(1)=w-w^{3}-w^{5}+w^{7}=(1-w^{2})(w-w^{5})
=w(1-w^{2})(1-w^{4})=2w(1-w^{2}),
\]
since $w^{8}=1$ and $w^{4}=-1$. Consequently,
\[
\tau(1)^{2}=4w^{2}(1-2w^{2}+w^{4})=-8w^{4}=8.
\]
As in Theorem \ref{thm:V.5.1} we show that
\[
\tau(1)^{p}=\tau(p)=\theta(p)\tau(1).
\]
Next we see that
\[
\theta(p)=(\tau(1)^{2})^{(p-1)/2}=8^{(p-1)/2}
=\left(\frac{8}{p}\right)\ \text{(Proposition \ref{prop:V.5.1})}
=\left(\frac{2}{p}\right)^{3}=\left(\frac{2}{p}\right).
\]
Therefore, $\left(\frac{2}{p}\right)=\theta(p)$. Now, by a direct
calculation involving $x=1,3,5,7$ (or, more easily, $x=1,3,-3,-1$), it can
be shown that $\theta(x)=(-1)^{(x^{2}-1)/8}$ and that
$(-1)^{(x^{2}-1)/8}$ depends only on the residue class of $x$ mod $8$.
\end{proof}

\begin{statement}{Example}
An application of the reciprocity law. The law of quadratic
reciprocity and the complementary formulas make it possible to calculate
the Legendre symbol by successive reductions. Let us calculate
$\left(\frac{23}{59}\right)$ without computing squares modulo $59$.
\[
\begin{aligned}
\left(\frac{23}{59}\right)
&=(-1)^{11\cdot29}\left(\frac{59}{23}\right)
=-\left(\frac{13}{23}\right)
=-(-1)^{6\cdot11}\left(\frac{23}{13}\right)
=-\left(\frac{10}{13}\right)\\
&=-\left(\frac{-3}{13}\right)
=-\left(\frac{-1}{13}\right)\left(\frac{3}{13}\right)
=-(-1)^{6}\left(\frac{3}{13}\right)\\
&=-(-1)^{6\cdot1}\left(\frac{13}{3}\right)
=-\left(\frac{1}{3}\right)=-1.
\end{aligned}
\]
Therefore, $23$ is not a square modulo $59$.
\end{statement}

%% ----------------------------------------------------------------------
\sechead{6}{The two squares theorem}{THE TWO SQUARES THEOREM}

We are going to apply Proposition \ref{prop:V.4.1} of \hyperref[sec:5.4]{§4} to the field
$L=\Q[i]$ where $i^{2}=-1$. Since $-1\equiv3\pmod4$ the ring $B$ of
integers of $L$ is $\Z+\Z i$. This ring is called the \emph{ring of
Gaussian integers}. Its discriminant is $-4$ (\hyperref[sec:5.3]{§3}, example). If $p$ is an
odd prime number and if $u$ is a generator of the cyclic group $\F_{p}^{*}$,
then $-1=u^{(p-1)/2}$. Therefore, $-1$ is a square in $\F_{p}$ if and only
if $(p-1)/2$ is even. We have the following classification for prime
numbers:
\begin{itemize}
\item $2$ ramifies in $\Q[i]$;
\item the primes of the form $4k+1$ split; and
\item the primes of the form $4k+3$ remain prime.
\end{itemize}

The following result will prove useful.

\begin{proposition}
\label{prop:V.6.1}
The ring $B=\Z+\Z i$ of Gaussian integers is principal.
\end{proposition}

\begin{proof}
We use the full force of an earlier result to prove this easy proposition.
With the notations of \hyperref[sec:4.3]{Chapter IV, §3}, we have $n=2$, $r_{1}=0$, $r_{2}=1$,
and $d=-4$. Therefore (\hyperref[sec:4.3]{Chapter IV, §3}, \hyperref[cor:IV.3.1]{the corollary of Proposition 1}),
\[
\textit{every ideal class of } B \textit{ contains an integral ideal of norm }
\leq\frac{4}{\pi}\frac{2}{4}|4|^{1/2}=\frac{4}{\pi}.
\]
Therefore, every ideal class contains the unit ideal, $B$ itself, which is
the only integral ideal of norm $1$ (note that $4/\pi<2$). Thus every ideal
of $B$ is equivalent to the principal ideal $B\cdot1$, so $B$ is a principal
ideal ring.
\end{proof}

\begin{statement}{A sketch of an elementary proof}
The points of $B$ may be identified with $\Z^{2}\subset\R^{2}=\C$. With the
usual identification of $\R^{2}$ with the plane (the square of Euclidean
distance is the norm in $\Q[i]$) the points of $\Z^{2}$ become the vertices
of squares covering the plane. A little geometry shows that, for any
$x\in\Q[i]$, there exists $z\in B$ such that
$\mathrm{N}(x-z)=|x-z|^{2}\leq(1/\sqrt{2})^{2}=\frac{1}{2}<1$. Now in
$\mathfrak{a}$, a non-zero ideal of $B$, choose a non-zero element $u$ of
minimum norm (NB. the norm is a positive integer-valued function on
$\Z[i]-(0)$). For any $v\in\mathfrak{a}$ there is a $z\in B$ such that
$\mathrm{N}(v/u-z)<1$. Therefore, $\mathrm{N}(v-zu)<\mathrm{N}(u)$, so
$v-zu=0$ (since $v-zu\in\mathfrak{a}$). Consequently,
$v\in Bu$ and $\mathfrak{a}=Bu$. The reader will observe the analogy with
the proof that $\Z$ is principal (\hyperref[sec:1.1]{Chapter I, §1}).
\end{statement}

\begin{proposition}
\label{prop:V.6.2}
(Fermat) Any prime number $p\equiv1\pmod4$ may be represented as the sum
of two squares (i.e.\ is of the form $p=a^{2}+b^{2}$ with $a,b\in\N$).
\end{proposition}

\begin{proof}
In fact $Bp=\mathfrak{p}_{1}\mathfrak{p}_{2}$, a product of distinct prime
ideals. Clearly $p^{2}=\mathrm{N}(Bp)
=\mathrm{N}(\mathfrak{p}_{1})\mathrm{N}(\mathfrak{p}_{2})$ (by \hyperref[sec:3.5]{Chapter III,
§5}, Proposition \ref{prop:III.5.2}). As the norms of $\mathfrak{p}_{1}$ and
$\mathfrak{p}_{2}$ are larger than $1$, necessarily
$\mathrm{N}(\mathfrak{p}_{1})=\mathrm{N}(\mathfrak{p}_{2})=p$. Now,
$\mathfrak{p}_{1}$ is a principal ideal $B\cdot(a+bi)$ ($a,b\in\Z$)
(Proposition \ref{prop:V.6.1}), so
$p=\mathrm{N}(a+bi)=a^{2}+b^{2}$.
\end{proof}

\begin{theorem}
\label{thm:V.6.1}
Let $x$ be a natural number and let
$x=\prod_{p}p^{\nu_{p}(x)}$ be its expression as a product of powers of
primes. For $x$ to be a sum of two squares it is necessary and sufficient
that, for every $p\equiv3\pmod4$, the exponent $\nu_{p}(x)$ be even.
\end{theorem}

\begin{proof}
In order to prove sufficiency we note that a sum of two squares
$a^{2}+b^{2}$ is the norm $\mathrm{N}(a+bi)$ of an element of $B$. By the
multiplicativity of norms, the set $S$ of sums of two squares is itself
stable under multiplication. Since $2=1^{2}+1^{2}\in S$ and since any
square is also an element of $S$ ($x^{2}=x^{2}+0^{2}$), sufficiency follows
from Proposition \ref{prop:V.6.2}.

Conversely let $x=a^{2}+b^{2}=(a+bi)(a-bi)$ be a sum of two squares
($a,b\in\N$) and let $p$ be a prime number $\equiv3\pmod4$. We have seen
that the ideal $Bp$ of $B$ is prime (Proposition \ref{prop:V.4.1}, \hyperref[sec:5.4]{§4}). Let
$n$ be the exponent of $Bp$ in the expression for $B\cdot(a+bi)$ as a
product of primes. As $Bp$ is stable under the automorphism
$\sigma:u+iv\mapsto u-iv$ of $B$ and since $\sigma(a+bi)=a-bi$, the
exponent of $Bp$ in the expression for $B\cdot(a-bi)$ is also $n$.
Therefore, in the expression for $B(a^{2}+b^{2})$, the exponent for $Bp$ is
$2n$. As no prime number besides $p$ belongs to $Bp$ (for
$Bp\cap\Z=p\Z$), we see that $\nu_{p}(x)=2n$ and $\nu_{p}(x)$ is even.
\end{proof}

%% ----------------------------------------------------------------------
\sechead{7}{The four squares theorem}{THE FOUR SQUARES THEOREM}

In this section we intend to prove the following theorem:

\begin{theorem}
\label{thm:V.7.1}
(Lagrange) Every natural number may be represented as the sum of four
squares.
\end{theorem}

The idea behind this proof is analogous to that of \hyperref[sec:5.6]{§6}: in place of the ring
of Gaussian integers, we will work in a suitably chosen ring of
\emph{quaternions}.

Let us begin with a definition of quaternions. Given a ring $A$, we will
write $(1,i,j,k)$ for the canonical base of the $A$-module $A^{4}$ and
define a multiplication on $A^{4}$ as follows:
\begin{equation}
\left\{
\begin{array}{l}
1 \text{ is the unit element,}\\
i^{2}=j^{2}=k^{2}=-1,\\
ij=-ji=k,\quad jk=-kj=i,\quad\text{and}\quad ki=-ik=j.
\end{array}
\right.
\label{eq:V.7.1}
\end{equation}
We extend this multiplication to elements $a+bi+cj+dk$ of $A^{4}$ by
linearity. Distributivity is then clear. It suffices to verify associativity
for elements of the base, e.g.\ $i(jk)=i^{2}=-1=k^{2}=(ij)k$. The formulas
in which $1$ appears being clear, there remain $3^{3}-1=26$ formulas to
check. The patient and incredulous reader will reduce the number of formulas
by the use of permutations and check those which remain. Others will accept
the author's claim that the given multiplication is associative. Provided
with this multiplication, $A^{4}$ is a not necessarily commutative ring, and
even an $A$-algebra, which we call the \emph{algebra of quaternions over}
$A$, denoted by $\mathbf{H}(A)$ (H in honor of W.~R. Hamilton, the inventor
of quaternions).

Given a quaternion $z=a+bi+cj+dk$ over $A$ (we write $a$ in place of
$a\cdot1$), we define the \emph{conjugate quaternion} of $z$ to be the
quaternion $\bar z=a-bi-cj-dk$.

\begin{lemma}
\label{lem:V.7.1}
For any $z$ and $z'\in\mathbf{H}(A)$,
$\overline{z+z'}=\bar z+\overline{z'}$,
$\overline{zz'}=\overline{z'}\bar z$, and $\overline{\bar z}=z$. In other
words, $z\mapsto\bar z$ is an involutive antiautomorphism of
$\mathbf{H}(A)$.
\end{lemma}

\begin{proof}
The first and the third formulas are clear. By linearity we see that to
prove the second it suffices to check that
$\overline{xy}=\bar y\bar x$ for $x,y\in\{1,i,j,k\}$. This is clear if $x=1$
or $y=1$. If $x=y=i$, then
\[
\overline{xy}=-1\quad\text{and}\quad\bar y\bar x=(-i)(-i)=i^{2}=-1.
\]
If $x=i$ and $y=j$, then
\[
\overline{xy}=\bar k=-k\quad\text{and}\quad
\bar y\bar x=(-j)(-i)=ji=-k.
\]
The others follow by a permutation argument.
\end{proof}

Given a quaternion $z$ over $A$, we call the quaternion
$\mathrm{N}(z)=z\bar z$ the \emph{reduced norm} of $z$.

\begin{lemma}
\label{lem:V.7.2}
\begin{alist}
\item For $z=a+bi+cj+dk\in\mathbf{H}(A)$,
$\mathrm{N}(z)=a^{2}+b^{2}+c^{2}+d^{2}$ (four squares!), an element of $A$.
\item For $z,z'\in\mathbf{H}(A)$,
$\mathrm{N}(zz')=\mathrm{N}(z)\mathrm{N}(z')$.
\end{alist}
\end{lemma}

\begin{proof}
To prove (a) we observe that, in the expansion of
$(a+bi+cj+dk)(a-bi-cj-dk)$, the cross-product terms cancel, leaving
$a^{2}+b^{2}+c^{2}+d^{2}$. To prove (b), observe that
\[
\mathrm{N}(zz')=zz'\overline{zz'}=zz'\bar z'\bar z
=z\mathrm{N}(z')\bar z=z\bar z\mathrm{N}(z')
\]
(since $\mathrm{N}(z')\in A$ and any element of $A$ commutes with every
quaternion), so $\mathrm{N}(zz')=\mathrm{N}(z)\mathrm{N}(z')$.
\end{proof}

Lemma \ref{lem:V.7.2} implies that, for a (commutative) ring $A$, the set of
reduced norms of quaternions, i.e.\ sums of four squares, is stable under
multiplication.

Now we shall study in greater detail the quaternion algebra
$\mathbf{H}(\Q)$, the non-commutative subring $\mathbf{H}(\Z)$, and the
subset $\mathbf{H}$ of ``quaternions of Hurwitz'', i.e.\ the quaternions of
the form $a+bi+cj+dk$, where either $a,b,c$, and $d$ belong to $\Z$ or all
four coefficients belong to $\frac{1}{2}+\Z$.

\begin{lemma}
\label{lem:V.7.3}
\begin{alist}
\item The set $\mathbf{H}$ of Hurwitz quaternions is a non-commutative
subring of $\mathbf{H}(\Q)$ containing $\mathbf{H}(\Z)$ and stable under
conjugation ($z\mapsto\bar z$).
\item For any $z\in\mathbf{H}$, $z+\bar z\in\Z$ and
$\mathrm{N}(z)=z\bar z\in\Z$.
\item In order that $z\in\mathbf{H}$ be invertible, it is necessary and
sufficient that $\mathrm{N}(z)=1$.
\item Every left ideal (respectively, right ideal) $\mathfrak{a}$ of
$\mathbf{H}$ is principal, i.e.\ of the form $\mathbf{H}z$ (respectively,
$z\mathbf{H}$).
\end{alist}
\end{lemma}

\begin{proof}
In (a) all assertions are clear except the stability of $\mathbf{H}$ under
multiplication. To complete the proof of (a) it suffices to show that, for
$u=\frac{1}{2}(1+i+j+k)$, $u\cdot1$, $u\cdot i$, $u\cdot j$, $u\cdot k$, and
$u^{2}\in\mathbf{H}$. We have
\[
u\cdot1=u,\quad u\cdot i=\tfrac{1}{2}(-1+i+j-k),\quad
u\cdot j=\tfrac{1}{2}(-1-i+j+k),
\]
\[
\text{and}\quad u\cdot k=\tfrac{1}{2}(-1+i-j+k),
\]
all elements of $\mathbf{H}$. By addition,
$2u^{2}=\frac{1}{2}(-2+2i+2j+2k)$, so $u^{2}\in\mathbf{H}$, too. To prove
(b) let
\[
z=\tfrac{1}{2}+a+(\tfrac{1}{2}+b)i+(\tfrac{1}{2}+c)j+(\tfrac{1}{2}+d)k,
\quad\text{with } a,b,c,d\in\Z.
\]
Then
\[
z+\bar z=1+2a\in\Z
\]
and
\[
z\bar z=(\tfrac{1}{2}+a)^{2}+(\tfrac{1}{2}+b)^{2}
+(\tfrac{1}{2}+c)^{2}+(\tfrac{1}{2}+d)^{2}
\in\tfrac{4}{4}+\Z\subset\Z.
\]
The preceding formula follows from Lemma \ref{lem:V.7.2}.

If $z$ is invertible in $\mathbf{H}$ and if $z'=z^{-1}$, then
\[
\mathrm{N}(z)\mathrm{N}(z')=\mathrm{N}(zz')=1.
\]
Since $\mathrm{N}(z)$ and $\mathrm{N}(z')$ are both positive integers ((b)
and Lemma \ref{lem:V.7.2}, (a)), $\mathrm{N}(z)=1$. Conversely, if
$z\in\mathbf{H}$ and if $\mathrm{N}(z)=1$, then
\[
z\bar z=\bar z z=\mathrm{N}(z)=1,
\]
so, since $\bar z\in\mathbf{H}$ by (a), $z$ is invertible in
$\mathbf{H}$. This proves (c).

To prove (d) take $x=a+bi+cj+dk\in\mathbf{H}(\Q)$. There exist
$a',b',c',d'\in\Z$ such that
\[
|a-a'|\leq\tfrac{1}{2},\quad |b-b'|\leq\tfrac{1}{2},\quad
|c-c'|\leq\tfrac{1}{2},\quad\text{and}\quad |d-d'|\leq\tfrac{1}{2}.
\]
Put $z=a'+b'i+c'j+d'k$. Then
\[
\mathrm{N}(x-z)=(a-a')^{2}+(b-b')^{2}+(c-c')^{2}+(d-d')^{2}
\leq4\cdot\tfrac{1}{4}=1.
\]
The inequality is even strict, except when $a,b,c$, and $d$ all belong to
$\frac{1}{2}+\Z$. But in this case $x\in\mathbf{H}$. Thus, for any quaternion
$x\in\mathbf{H}(\Q)$, there exists a Hurwitz quaternion
$z\in\mathbf{H}$ such that $\mathrm{N}(x-z)<1$ (it is precisely because we
must have strict inequality that we work with the Hurwitz quaternions;
$\mathbf{H}(\Z)$ would not be sufficient). Now let $\mathfrak{a}$ be a left
ideal of $\mathbf{H}$. To show that $\mathfrak{a}$ is principal we may
assume $\mathfrak{a}\neq(0)$. Choose in $\mathfrak{a}$ a non-zero element
$u$ for which $\mathrm{N}(u)$ is a minimum (such $u$ exists, since the
reduced norm is a positive integer-valued function on
$\mathbf{H}-(0)$, by (b)). Clearly, $u$ is invertible in $\mathbf{H}(\Q)$
with inverse $\bar u\mathrm{N}(u)^{-1}$ (this shows that $\mathbf{H}(\Q)$ is
a skew field). For $y\in\mathfrak{a}$ form
$yu^{-1}\in\mathbf{H}(\Q)$ and take an element $z\in\mathbf{H}$ such that
$\mathrm{N}(yu^{-1}-z)<1$. Then, by Lemma \ref{lem:V.7.2}, (b), we have
\[
\mathrm{N}(y-zu)=\mathrm{N}((yu^{-1}-z)u)<\mathrm{N}(u).
\]
Since $y-zu\in\mathfrak{a}$ and since $\mathrm{N}(u)$ is as small as
possible, it follows that $y-zu=0$, $y\in\mathbf{H}u$, and
$\mathfrak{a}=\mathbf{H}u$.
\end{proof}

Now we are ready to prove Theorem \ref{thm:V.7.1}. Since the set of elements
of $\Z$ which are sums of four squares is multiplicatively stable (cf.\
Lemma \ref{lem:V.7.2}), Theorem \ref{thm:V.7.1} follows from the following
proposition.

\begin{proposition}
\label{prop:V.7.1}
Any prime number is the sum of four squares.
\end{proposition}

\begin{proof}
Since $2=1^{2}+1^{2}+0^{2}+0^{2}$, we may suppose that $p$ is odd. Since $p$
commutes with any quaternion, the left ideal $\mathbf{H}p$ is two-sided.
Consider the quotient ring $\mathbf{H}/\mathbf{H}p$. Since $p$ is odd, any
$z\in\mathbf{H}$ is congruent mod $\mathbf{H}p$ to an element of
$\mathbf{H}(\Z)$ (if the components of $z$ all belong to
$\frac{1}{2}+\Z$, form $z+p\cdot\frac{1}{2}(1+i+j+k)$). Therefore,
$\mathbf{H}/\mathbf{H}p$ is isomorphic to the corresponding quotient ring of
$\mathbf{H}(\Z)$, i.e.\ $\mathbf{H}(\F_{p})$.

Since the quadratic form $a^{2}+b^{2}+c^{2}+d^{2}$ represents $0$ in
$\F_{p}$ (\hyperref[sec:1.7]{Chapter I, §7}, Theorem \ref{thm:I.7.2}; see the remark below for a
direct proof), $\mathbf{H}(\F_{p})$ contains a non-zero element of reduced
norm zero. Such an element is not invertible (Lemma \ref{lem:V.7.2}, (b)),
so it generates a non-trivial left ideal. It follows that $\mathbf{H}p$ is
properly contained in a left ideal $\mathbf{H}z$ distinct from
$\mathbf{H}$. Consequently, $p=z'z$ for some $z,z'\in\mathbf{H}$ both
non-units. Thus,
\[
p^{2}=\mathrm{N}(p)=\mathrm{N}(z)\mathrm{N}(z')
\]
and, since $\mathrm{N}(z)$ and $\mathrm{N}(z')$ are both integers strictly
larger than one (Lemma \ref{lem:V.7.3}, (b) and (c)),
$\mathrm{N}(z)=\mathrm{N}(z')=p$.

Put $z=a+bi+cj+dk$ ($a,b,c,d\in\Z$ or in $\frac{1}{2}+\Z$). If
$a,b,c,d\in\Z$, then $p=\mathrm{N}(z)=a^{2}+b^{2}+c^{2}+d^{2}$ and we are
finished. It suffices to show that, if $a,b,c,d\in\frac{1}{2}+\Z$. We may
make a reduction to the preceding case by multiplying $z$ by an element of
$\mathbf{H}$ of reduced norm $1$, more precisely by an element of the form
$\frac{1}{2}(\pm1\pm i\pm j\pm k)$. To see this consider the residue class
$\eta$ associated with $2z$ in
$\mathbf{H}(\Z)/4\mathbf{H}(\Z)\cong\mathbf{H}(\Z/4\Z)$. Since
$\mathrm{N}(z)\in\Z$, $\mathrm{N}(2z)\in4\Z$, so $\mathrm{N}(\eta)=0$ and
$\eta\bar\eta=0$. We observe that $\bar\eta$ is the residue class of a
quaternion $z'$ of the form $\pm1\pm i\pm j\pm k$. Set
$u=\frac{1}{2}z'\in\mathbf{H}$; then $u$ is of reduced norm $1$ and,
inasmuch as $(2z)(2u)$ is zero mod $4$, $zu\in\mathbf{H}(\Z)$. Now
$p=\mathrm{N}(z)=\mathrm{N}(zu)$, so we are finished.
\end{proof}

\begin{statement}{Remark}
Here is a very elementary proof of the fact that, over a finite field $K$,
the quadratic form $a^{2}+b^{2}+c^{2}+d^{2}$ represents $0$ (i.e.\ has a
non-trivial zero in $K^{4}$). It suffices to take $c=1$ and $d=0$ and to
show that the equation $a^{2}+b^{2}+1=0$ has a solution in $K^{2}$. Write
this equation in the form $b^{2}+1=-a^{2}$. In characteristic $2$, we may
take $b=0$ and $a=1$. If $\operatorname{card}(K)=q$ is odd, there are
$(q+1)/2$ squares in $K$ ($0$ and the $(q-1)/2$ non-zero squares). Thus the
set $T$ (respectively, $T'$) of elements of $K$ of the form $b^{2}+1$ with
$b\in K$ (respectively, of the form $-a^{2}$ with $a\in K$) consists of
$(q+1)/2$ elements. Since $(q+1)/2+(q+1)/2>q$, $T\cap T'\neq\phi$, so
$b^{2}+1=-a^{2}$ has a solution.
\end{statement}
